% Einheit 1 von 4 – Normalparabel und Streckfaktor (bank/quadratische-funktionen/e1.jsonl, zusammenbau v0.1)
%% TODO Zweigzeile Teil 1: was hier gelernt wird (Satz in Schülersprache aus dem Einheitstitel) – Einheit 1
\einheitenkopf[e1]{Einheit 1 von 4 · Normalparabel und Streckfaktor}
\zweigzeile{ab Kl. 9, je nach Buch bis Kl. 10 · P10 oft}
\setcounter{aufgabe}{10}

%% TODO Titel: Ich-kann-Satz fehlt in der Bank (Platzhalter: Kettenname) – Nr. 11
%% TODO Anweisung: ein Satz über der ersten Teilaufgabe fehlt in der Bank – Nr. 11
\begin{aufgabe}{Normalparabel und Streckfaktor}
\begin{teile}
\teil $f(x) = 4x - 1$ – Gerade oder Parabel? \kreuz{Gerade} \kreuz{Parabel}
\end{teile}
%% TODO Darstellung für schwach fehlt in der Bank (quadratische-funktionen-e1-k1-s1-v1; Abschnitt „Für schwache Schüler“)
\swz{Wertetabelle zu $f(x) = x^2$ – fülle aus.
\wertetabelle{x}{f(x)}{-2,-1,0,1,2}}{}
%% TODO Darstellung für schwach fehlt in der Bank (quadratische-funktionen-e1-k1-s1-v2; Abschnitt „Für schwache Schüler“)
\swz{Wertetabelle zu $f(x) = x^2$ – fülle aus.
\wertetabelle{x}{f(x)}{-3,-2,-1,0,1,2,3}}{}
%% TODO Darstellung für schwach fehlt in der Bank (quadratische-funktionen-e1-k1-s1-v3; Abschnitt „Für schwache Schüler“)
\swz{Wertetabelle zu $f(x) = x^2$ – fülle aus.
\wertetabelle{x}{f(x)}{-4,-2,0,2,4}}{}
%% TODO Darstellung für schwach fehlt in der Bank (quadratische-funktionen-e1-k1-s1-v4; Abschnitt „Für schwache Schüler“)
\swz{Wertetabelle zu $f(x) = x^2$ – fülle aus.
\wertetabelle{x}{f(x)}{-5,-3,-1,1,3}}{}
%% TODO Darstellung für schwach fehlt in der Bank (quadratische-funktionen-e1-k1-s1-v5; Abschnitt „Für schwache Schüler“)
\swz{Wertetabelle zu $f(x) = x^2$ – fülle aus.
\wertetabelle{x}{f(x)}{-6,-4,-1,3,5}}{}
\swfrage{Was bleibt gleich, was ändert sich?}
\swa{Trage die Punkte ein und zeichne die Parabel zu $f(x) = x^2$ als Bogen.
\wertetabelle[9,4,1,0,1,4,9]{x}{f(x)}{-3,-2,-1,0,1,2,3}}{\begin{ksys}[xmin=-4,xmax=4,ymin=-1,ymax=10]
\end{ksys}}
\swz{Normalparabel $f(x) = x^2$ – Scheitel und Symmetrieachse?}{\begin{ksys}[xmin=-4,xmax=4,ymin=-1,ymax=9,ablesen]
\parabel{1}{0}{0}{f}
\end{ksys}}
%% TODO Darstellung für schwach fehlt in der Bank (quadratische-funktionen-e1-k1-s4-v1; Abschnitt „Für schwache Schüler“)
\swz{$f(x) = x^2$ – $f(-6)$?}{}
%% TODO Darstellung für schwach fehlt in der Bank (quadratische-funktionen-e1-k1-s5-v1; Abschnitt „Für schwache Schüler“)
\swz{Wertetabelle zu $f(x) = 4x^2$ – fülle aus.
\wertetabelle{x}{f(x)}{-2,-1,0,1,2}}{}
\begin{teile}
\teil $f(x) = -4x^2$ – \kreuz{nach oben} \kreuz{nach unten} \\ \kreuz{schmaler} \kreuz{breiter} als die Normalparabel
\end{teile}
\begin{teile}
\teil Welche Tabelle gehört zu $f(x) = 0{,}5x^2$? \\ A: \wertetabelle[2,0{,}5,0,0{,}5,2]{x}{y}{-2,-1,0,1,2} \\ B: \wertetabelle[1,0{,}25,0,0{,}25,1]{x}{y}{-2,-1,0,1,2} \\ C: \wertetabelle[4,1,0,1,4]{x}{y}{-2,-1,0,1,2} \\ \kreuz{A} \kreuz{B} \kreuz{C}
\end{teile}
\swa{$f(x) = 0{,}5x^2$ – zeichne die Parabel mit den halbierten Werten der Normalparabel.}{\begin{ksys}[xmin=-5,xmax=5,ymin=-1,ymax=9]
\end{ksys}}
\begin{teile}
\teil Welche Gleichung gehört zum Graphen? \\ \kreuz{$f(x) = x^2 + 3$} \\ \kreuz{$f(x) = -x^2 - 3$} \\ \kreuz{$f(x) = x^2 - 3$} \\ \kreuz{$f(x) = -x^2 + 3$}
\end{teile}
\begin{teile}
\teil Welche Wertetabelle gehört zu $y = 4x^2$? \\ A: \wertetabelle[-8,-4,0,4,8]{x}{y}{-2,-1,0,1,2} \\ B: \wertetabelle[16,4,0,4,16]{x}{y}{-2,-1,0,1,2} \\ C: \wertetabelle[64,16,0,16,64]{x}{y}{-2,-1,0,1,2} \\ \kreuz{A} \kreuz{B} \kreuz{C} \hfill (P10 2026 FOR)
\end{teile}
\begin{teile}
\teil Ein Stein fällt von einer $45\,$m hohen Brücke. Der Graph zeigt seine Höhe $h$ (in m) nach $t$ Sekunden. Welche Gleichung passt? Kreuze an und begründe. \\ \kreuz{$h(t) = 5t^2 + 45$} \\ \kreuz{$h(t) = 5t^2 - 45$} \\ \kreuz{$h(t) = -5t^2 + 45$} \\ \kreuz{$h(t) = -5t^2 - 45$} \hfill (P10 2015 OS)
\end{teile}
\end{aufgabe}

%% TODO Titel: Ich-kann-Satz fehlt in der Bank (Platzhalter: Kettenname) – Nr. 12
%% TODO Anweisung: ein Satz über der ersten Teilaufgabe fehlt in der Bank – Nr. 12
\begin{aufgabe}{Normalparabel und Streckfaktor – Fehler finden, Begründen, Darstellungswechsel, Anwendung}
\swz[3]{Mia füllt die Tabelle zu $f(x) = x^2$ so aus: \\ \wertetabelle[-16,-4,0,4,16]{x}{f(x)}{-4,-2,0,2,4} \\ Finde den Fehler und korrigiere.}{}
\swfrage{$f(x) = x^2$ hat bei $x = 8$ und bei $x = -8$ denselben Wert. Warum?}
\swz{Die Tabelle gehört zu $f(x) = a \cdot x^2$. \\ \wertetabelle[7,28,63]{x}{f(x)}{1,2,3} \\ Wie heißt $a$?}{}
\swz[3]{Der Bogen einer Brücke folgt $h(x) = -0{,}02x^2$ ($x$ und $h$ in m, der Ursprung liegt im höchsten Punkt des Bogens). Wie tief liegt der Bogen $20\,$m neben der Mitte unter dem höchsten Punkt?}{\begin{ksys}[xmin=-30,xmax=30,ymin=-20,ymax=5,xstep=10,ystep=5,xlabel=$x$ in m,ylabel=$h$ in m]
\funktion{-0.02*\x^2}{h}
\end{ksys}}
\end{aufgabe}

\uebersichtskasten{Normalparabel: p(x) = x². Scheitel S(0 | 0), Symmetrieachse ist die y-Achse, nach oben geöffnet, kein y-Wert ist negativ. \\ Vom Scheitel aus: 1 nach rechts 1 nach oben, 2 nach rechts 4 nach oben, 3 nach rechts 9 nach oben – nach links genauso, denn (−3)² = 9. \\ Streckfaktor: p(x) = a·x². a negativ: nach unten geöffnet (−x² ist das Spiegelbild). a zwischen 0 und 1: breiter (gestaucht). a größer als 1: schmaler (gestreckt). \\ p(x) = 2x²: jeder y-Wert doppelt so groß, Punkte (1 | 2) und (2 | 8).}
